\section{Introduction}
\label{sec:intro}

Text-to-Image (T2I) diffusion models~\cite{sd,wuerstchen,ddpm} have demonstrated extraordinary capabilities in synthesizing visually compelling and photorealistic imagery from natural language descriptions. These foundation models are increasingly harnessed across diverse creative industries, including digital art creation, video game asset design, personalized concept synthesis, and image editing~\cite{prompt2prompt,instructpix2pix}. Content creators and digital artists typically demand granular, independent control over both the semantic content and the visual style of synthesized images to fulfill their creative vision. While the semantic content of a scene (e.g., "a vintage clock on a wooden table" or "a sitting cat") can be precisely articulated via descriptive text prompts, conveying an artist's signature visual style---characterized by intricate brushstroke impasto, delicate pigment dispersion, non-Euclidean geometric tiling, and distinctive physical medium properties (e.g., oil painting, stained glass, woodblock printing, origami, glowing neon cyberpunk)---is inherently non-verbal, abstract, and vastly more complex. Consequently, communicating these subtle stylistic nuances through text prompts alone is severely hindered by prompt ambiguity and lexical constraints. This has established visual style personalization via visual prompting from a single exemplar image as an indispensable frontier in generative computer vision~\cite{stylealigned,instantstyle,rout2025rbmodulation}.

% ==============================================================================
% Teaser Figure: Comprehensive Qualitative Matrix (Content x Style)
% ==============================================================================
\begin{figure*}[t]
\centering
\includegraphics[width=\textwidth]{figs/qualitative_matrix_teaser.png}
\caption{\textbf{Qualitative Stylization Matrix of \method across Diverse Content Categories and Artistic Styles.} 
Rows depict representative content targets spanning diverse semantic domains (animals, household artifacts, and footwear); columns correspond to diverse artistic styles (oil paintings, mosaic tile art, woodblock prints, glowing neon cyberpunk, and 3D renders). All results are synthesized under null-prompt inference (without text prompt conditioning, akin to RB-Modulation~\cite{rout2025rbmodulation}), demonstrating faithful subject geometry retention while authentically rendering rich stroke dynamics and textural nuances without semantic content leakage.}
\label{fig:teaser_matrix}
\end{figure*}

To personalize pre-trained T2I diffusion models for a target artistic style, existing research predominantly follows two overarching paradigms: (1) fine-tuning-based optimization and (2) training-free inference-time modulation. Fine-tuning approaches---exemplified by DreamBooth~\cite{dreambooth}, Custom Diffusion, LoRA~\cite{lora}, ZipLoRA~\cite{ziplora}, StyleDrop~\cite{styledrop}, and B-LoRA~\cite{b-lora}---optimize the model's text embeddings, cross-attention projection matrices, or low-rank residual weights via the denoising diffusion objective. When fine-tuned on a curated set of multiple reference images (typically 4--10 exemplars of the same artistic domain), these methods can effectively capture stylistic patterns. However, they suffer from three fundamental drawbacks:
(i) curating multiple high-quality images sharing an identical artistic style is often impractical or impossible, particularly for unique historical artworks or bespoke digital designs;
(ii) optimizing model weights for every novel style incurs prohibitive computational cost and latency (requiring several minutes of GPU gradient descent per concept); and
(iii) critically, when constrained to a \textit{single reference exemplar}, fine-tuning methods suffer from severe semantic memorization and catastrophic overfitting---the network inadvertently memorizes and reproduces the concrete semantic objects depicted in the reference artwork (e.g., a specific portrait or building) rather than isolating its abstract stylistic essence.

To overcome the heavy computational burden and multi-image dependency of fine-tuning, recent efforts have focused on training-free personalization by directly manipulating keys, values, and feature maps within the attention layers of pre-trained diffusion models at inference time~\cite{stylealigned,instantstyle,styleid,diffstyler,freestyle}. Nevertheless, existing training-free frameworks confront severe challenges in isolating style from semantic content and transferring it faithfully without structural artifacts. For instance, StyleAligned~\cite{stylealigned} and Swapping Self-Attention (SSA) synchronize style across images by sharing self-attention keys and values. However, these methods rely on deterministic DDIM inversion~\cite{DDIM} to cache reference representations, which inherently compromises high-frequency textural details and accumulates discretization errors; furthermore, they require dual coupled reverse diffusion processes that double GPU memory footprint and runtime. To circumvent DDIM inversion, InstantStyle~\cite{instantstyle} injects reference style features into heuristic subsets of cross-attention layers within a pre-trained IP-Adapter~\cite{ipadapter}. However, manually identifying optimal injection layers is brittle, ad-hoc, and fails to generalize across diverse network architectures. More detrimentally, linear feature injection directly induces catastrophic \textbf{content leakage}---inadvertently synthesizing foreign semantic objects from the style reference into the generated scene (e.g., painting extraneous dragons or houses when only a metallic golden texture was desired). To suppress geometric collapse, InstantStyle must rely on an external ControlNet~\cite{controlNet}, which severely restricts the generation to rigid spatial poses and stifles artistic diversity.

Most recently, RB-Modulation~\cite{rout2025rbmodulation} introduced a training-free framework leveraging stochastic optimal control with a terminal style cost, eliminating the dependency on external adapters. However, to prevent content leakage, RB-Modulation resorts to aggressive \textbf{mean-pooling} ($\bar{\mathbf{t}}_{\text{style}} = \frac{1}{L} \sum_{j=1}^L \mathbf{t}_{\text{style}}^{(j)}$) across all reference style tokens. While collapsing style tokens to a single spatial mean suppresses semantic object transfer, it catastrophically eliminates spatial token variance, resulting in acute \textbf{under-stylization}, muted brushstroke textures, and washed-out artistic fidelity. Furthermore, RB-Modulation's unconstrained gradient guidance updates push sampling trajectories off the valid diffusion data manifold, and its reliance on the supervised Consistent Style Descriptor (CSD) frequently triggers domain collapse on complex non-photorealistic artistic media. Consequently, existing methods remain locked in an acute polarization: either suffering from severe content collapse and semantic leakage (e.g., InstantStyle, DiffuseIT~\cite{diffuseIT}), or languishing in under-stylization and textural dilution (e.g., RB-Modulation, ITOC).

To decisively resolve this dilemma, we present \textbf{\method}, an end-to-end training-free style personalization framework built on cascaded diffusion architectures (Stable Cascade~\cite{wuerstchen}). \method establishes a new Pareto frontier in visual style transfer without test-time optimization, DDIM inversion, or text prompt conditioning. Rather than compromising between content fidelity and stylistic intensity, \method reconciles them through three complementary, mathematically principled mechanisms:
\begin{itemize}
    \item \textbf{Blended Style Tokens with Null-Space Orthogonal Projection}: We formulate a soft blending operator ($\alpha_s = 0.85$) that preserves high-frequency artistic texture variance while smoothing out isolated semantic anomalies. By mathematically projecting style tokens onto the orthogonal complement (null space) of content representations $\Vortho$ inside cross-attention layers, \method eliminates constructive feature interference and guarantees zero semantic leakage while restoring rich, dynamic textural expression.
    \item \textbf{Score-Orthogonal DINO Latent Guidance}: We introduce a terminal guidance loss formulated in the feature space of self-supervised Vision Transformers (DINO-ViT-S/16)~\cite{dino}. Because DINO is trained without discrete categorical labels, its representations naturally capture mid-level spatial co-occurrences, texture continuity, and stroke dynamics without semantic bias. By projecting the guidance gradients onto the null space of the reverse diffusion score vector, \method enforces manifold preservation, preventing out-of-distribution drift and eliminating oversaturation artifacts.
    \item \textbf{Early-Step AdaIN Pushforward Controller}: Anchored by the fundamental stochastic noise trade-offs established in SDEdit~\cite{sdedit}---wherein early timesteps govern coarse geometric layout while late steps refine fine textures---\method applies Adaptive Instance Normalization (AdaIN) directly to the predicted clean latent $\bx_0$ during the initial sampling steps. This rapidly anchors the global chromatic palette and contrast without corrupting late-stage brushstroke synthesis or perturbing the diffusion variance schedule.
\end{itemize}

Our primary contributions are summarized as follows:
\begin{enumerate}
    \item We propose \textbf{\method}, a plug-and-play, training-free framework for exemplar-based visual style transfer that establishes a new Pareto frontier without requiring network fine-tuning, DDIM inversion, or text prompt engineering.
    \item We introduce \textbf{Blended Style Tokens with Null-Space Orthogonal Projection and Energy Compensation}, which mathematically resolves the trade-off between the catastrophic content leakage of InstantStyle and the textural collapse of RB-Modulation's mean pooling.
    \item We formulate \textbf{Score-Orthogonal DINO Latent Guidance} with rigorous mathematical proofs of manifold stability, leveraging self-supervised ViT patch representations as an authentic texture harmonizer to align fine brushstrokes without semantic bias.
    \item We design an \textbf{Early-Step AdaIN Pushforward Controller} informed by SDEdit's stochastic trade-offs, coupled with semantic gated edge extraction, demonstrating decisive quantitative and qualitative superiority over six SOTA baselines on the comprehensive 225-pair SOICT-Data benchmark.
\end{enumerate}
